\section{语音模型（Omni）：听与说的统一}

\subsection{Thinker-Talker 架构}

Qwen-Omni 采用了一套 ``Thinker + Talker'' 架构：

\begin{itemize}
    \item \textbf{Thinker}：负责理解和推理
    \item \textbf{Talker}：负责语音合成和输出
    \item 实现了 end-to-end 的语音交互（不依赖 ASR 管线）
\end{itemize}

\begin{knowledgebox}{端到端语音 vs ASR 管线}
传统做法是 ASR → LLM → TTS 三段式管线，但 Qwen-Omni 走 end-to-end 路线——模型直接理解语音输入并生成语音输出。这不仅减少了延迟，还能捕捉语调、情感等非文字信息。TTS 方面可通过 prompt 描述声音特征来定制音色。
\end{knowledgebox}

当前 Omni 模型的语言智力略有降质（约 Gemini 2.5 Flash 水平），但语音能力对标 Gemini 2.5 Pro。

\subsection{本章小结}

语音是 Agent 与人类交互的重要通道。Qwen-Omni 的 Thinker-Talker 架构是实现``三进三出''（文本、视觉、音频的理解与生成）目标的关键一步。

